\documentclass[11pt,a4paper]{article}
\usepackage[utf8]{inputenc}
\usepackage[T1]{fontenc}
\usepackage{geometry}
\usepackage{multicol}
\usepackage[table]{xcolor}
\usepackage{colortbl}
\usepackage{listings}
\usepackage{tcolorbox}
\usepackage{enumitem}
\usepackage{titlesec}
\usepackage{fancyhdr}   
\usepackage{booktabs}
\usepackage{array}
\usepackage{tabularx}
\usepackage{fontawesome5}
\usepackage{hyperref}
\usepackage{tikz}
\usepackage{microtype}
\usepackage{iftex}
% Font fallbacks: keep the PDF building even if optional fonts are missing.
% Ensure sans-serif consistency
\IfFileExists{raleway.sty}{
    \usepackage[sfdefault]{raleway}
}{
    \usepackage{helvet}
    \renewcommand{\familydefault}{\sfdefault}
}
\IfFileExists{sourcecodepro.sty}{\usepackage{sourcecodepro}}{\usepackage{inconsolata}}
\usepackage{calc}
\usepackage{graphicx}
\tcbuselibrary{skins,breakable}

% Page geometry - optimized margins
\geometry{
    left=1cm,
    right=1cm,
    top=1.5cm,
    bottom=1cm
}

% Modern Color Palette - Dark theme inspired
\definecolor{bgdark}{RGB}{30, 34, 42}
\definecolor{bgcard}{RGB}{40, 44, 55}
\definecolor{primaryblue}{RGB}{86, 156, 214}
\definecolor{secondaryblue}{RGB}{78, 201, 176}
\definecolor{accentpurple}{RGB}{198, 120, 221}
\definecolor{accentorange}{RGB}{229, 192, 123}
\definecolor{textlight}{RGB}{220, 223, 228}
\definecolor{textmuted}{RGB}{150, 155, 165}
\definecolor{codegreen}{RGB}{152, 195, 121}
\definecolor{codered}{RGB}{224, 108, 117}
\definecolor{codeyellow}{RGB}{229, 192, 123}
\definecolor{warnorange}{RGB}{255, 170, 100}
\definecolor{successgreen}{RGB}{152, 195, 121}
\definecolor{dangerred}{RGB}{224, 108, 117}
\definecolor{infocyan}{RGB}{97, 175, 239}
\definecolor{codebg}{RGB}{50, 55, 68}
\definecolor{sectionline}{RGB}{86, 156, 214}

% Set page background
\pagecolor{bgdark}
\color{textlight}

% Hyperref setup
\hypersetup{
    colorlinks=true,
    linkcolor=primaryblue,
    urlcolor=secondaryblue
}

% Layout sanity: avoid overfull boxes in tight 4-column layout
\setlength{\parindent}{0pt}
\setlength{\parskip}{0pt}
\sloppy
\emergencystretch=1em

% Code listing style - modern dark theme
\lstdefinestyle{cstyle}{
    language=C,
    basicstyle=\ttfamily\fontsize{8.5pt}{10pt}\selectfont\color{textlight},
    keywordstyle=\color{accentpurple}\bfseries,
    stringstyle=\color{codegreen},
    commentstyle=\color{textmuted}\itshape,
    identifierstyle=\color{textlight},
    showstringspaces=false,
    breaklines=true,
    breakatwhitespace=false,
    columns=fullflexible,
    keepspaces=true,
    tabsize=2,
    frame=none,
    backgroundcolor=\color{codebg},
    xleftmargin=6pt,
    xrightmargin=6pt,
    aboveskip=6pt,
    belowskip=6pt,
    framesep=6pt,
    rulecolor=\color{bgcard},
    morekeywords={pthread_t, pthread_mutex_t, pthread_cond_t, pthread_attr_t, 
                  pthread_rwlock_t, pthread_barrier_t, pthread_spinlock_t,
                  pthread_key_t, pthread_once_t, size_t, sem_t, bool, NULL},
    emph={pthread_create, pthread_join, pthread_exit, pthread_self,
          pthread_mutex_lock, pthread_mutex_unlock, pthread_mutex_init,
          pthread_cond_wait, pthread_cond_signal, pthread_cond_broadcast,
          pthread_rwlock_rdlock, pthread_rwlock_wrlock, pthread_barrier_wait,
          sem_wait, sem_post, printf, malloc, free},
    emphstyle=\color{primaryblue},
    literate={*}{{\textcolor{accentorange}{*}}}1
             {&}{{\textcolor{accentorange}{\&}}}1
}
\lstset{style=cstyle}

% Section formatting - sleek modern look
\titleformat{\section}
    {\normalfont\fontsize{14pt}{16pt}\selectfont\bfseries\color{primaryblue}}
    {\thesection}{0em}{}
\titlespacing{\section}{0pt}{8pt}{4pt}

\titleformat{\subsection}
    {\normalfont\fontsize{11pt}{13pt}\selectfont\bfseries\color{secondaryblue}}
    {\thesubsection}{0em}{}
\titlespacing{\subsection}{0pt}{6pt}{3pt}

% Custom section command with icon and line
\newcommand{\styledSection}[2]{%
    \vspace{8pt}%
    \noindent\textcolor{primaryblue}{\rule{\linewidth}{1.5pt}}\\[-8pt]
    {\fontsize{14pt}{16pt}\selectfont\bfseries\color{primaryblue}#1\ #2}\\[-4pt]
    \textcolor{primaryblue!50}{\rule{\linewidth}{0.8pt}}%
    \vspace{4pt}%
}

\newcommand{\styledSubSection}[1]{%
    \vspace{4pt}%
    {\fontsize{11pt}{13pt}\selectfont\bfseries\color{secondaryblue}#1}%
    \vspace{2pt}%
}

% Custom boxes - modern card style
\tcbset{
    boxrule=0pt,
    arc=4pt,
    outer arc=4pt,
    left=6pt,
    right=6pt,
    top=5pt,
    bottom=5pt,
    fonttitle=\bfseries\fontsize{9pt}{11pt}\selectfont,
    before skip=6pt,
    after skip=6pt,
    sharp corners=south
}

\newtcolorbox{funcbox}[1][]{
    enhanced,
    colback=bgcard,
    colframe=primaryblue,
    boxrule=1pt,
    title=#1,
    coltitle=white,
    colbacktitle=primaryblue,
    attach boxed title to top left={yshift=-2mm, xshift=3mm},
    boxed title style={arc=2pt, outer arc=2pt}
}

\newtcolorbox{warnbox}{
    enhanced,
    colback=warnorange!8!bgdark,
    colframe=warnorange,
    boxrule=0pt,
    borderline west={3pt}{0pt}{warnorange},
    title={\footnotesize\faExclamationTriangle\ \textbf{Warning}},
    coltitle=warnorange,
    colbacktitle=warnorange!8!bgdark,
    fonttitle=\fontsize{9pt}{11pt}\selectfont\bfseries,
    attach boxed title to top left={yshift=-1mm, xshift=2mm},
    boxed title style={boxrule=0pt, arc=0pt}
}

\newtcolorbox{tipbox}{
    enhanced,
    colback=successgreen!8!bgdark,
    colframe=successgreen,
    boxrule=0pt,
    borderline west={3pt}{0pt}{successgreen},
    title={\footnotesize\faLightbulb\ \textbf{Tip}},
    coltitle=successgreen,
    colbacktitle=successgreen!8!bgdark,
    fonttitle=\fontsize{9pt}{11pt}\selectfont\bfseries,
    attach boxed title to top left={yshift=-1mm, xshift=2mm},
    boxed title style={boxrule=0pt, arc=0pt}
}

\newtcolorbox{infobox}{
    enhanced,
    colback=infocyan!8!bgdark,
    colframe=infocyan,
    boxrule=0pt,
    borderline west={2.5pt}{0pt}{infocyan},
    arc=3pt
}

\newtcolorbox{codebox}{
    enhanced,
    colback=codebg,
    colframe=bgcard,
    boxrule=0.5pt,
    arc=3pt,
    left=2pt,
    right=2pt,
    top=2pt,
    bottom=2pt
}

% Compact lists with custom bullets
\setlist{nosep, leftmargin=10pt, labelindent=0pt}
\setlist[itemize]{label=\textcolor{primaryblue}{\tiny\faChevronRight}}
\setlist[enumerate]{label=\textcolor{primaryblue}{\footnotesize\arabic*.}}

% Remove page numbers
\pagestyle{empty}

% Column separator - subtle gradient effect
\setlength{\columnseprule}{0.5pt}
\def\columnseprulecolor{\color{primaryblue!40}}
\setlength{\columnsep}{8pt}
\setlength{\multicolsep}{2pt}
\setlength{\premulticols}{0pt}
\setlength{\postmulticols}{0pt}

% Custom table style
\newcommand{\tablestyle}{%
    \renewcommand{\arraystretch}{1.3}%
    \fontsize{9pt}{11pt}\selectfont%
    \rowcolors{1}{bgcard}{bgdark}%
}

% Inline code style
\newcommand{\code}[1]{\texttt{\small\color{accentorange}#1}}

% Function name style
\newcommand{\func}[1]{\texttt{\color{primaryblue}#1}}

\begin{document}

% Title - Modern header design
\begin{tikzpicture}[remember picture, overlay]
    \fill[primaryblue!18!bgdark] (current page.north west) rectangle ([yshift=-0.95cm]current page.north east);
    \draw[primaryblue, line width=1.2pt] ([yshift=-0.95cm]current page.north west) -- ([yshift=-0.95cm]current page.north east);
\end{tikzpicture}

\begin{center}
    {\fontsize{24pt}{28pt}\selectfont\bfseries\color{primaryblue} \faCodeBranch\ POSIX Threads (pthreads) Cheat Sheet}\\[6pt]
    {\fontsize{11pt}{13pt}\selectfont\color{textmuted} Complete Reference for Multithreaded Programming in C \quad\textbar\quad \texttt{\#include <pthread.h>} \quad\textbar\quad Compile: \texttt{gcc -pthread}}
\end{center}

\vspace{8pt}

\begin{multicols}{2}
\raggedcolumns
\fontsize{9.5pt}{11pt}\selectfont

% ============================================================================
% DATA TYPES
% ============================================================================
\styledSection{\faDatabase}{Data Types}

\begin{tcolorbox}[enhanced, colback=bgcard, colframe=bgcard, boxrule=0pt, arc=3pt, left=3pt, right=3pt, top=2pt, bottom=2pt]
\tablestyle
\begin{tabular}{@{}>{\ttfamily\color{accentorange}}l@{\hspace{4pt}}>{\color{textmuted}}l@{}}
pthread\_t & Thread identifier \\
pthread\_attr\_t & Thread attributes \\
pthread\_mutex\_t & Mutex lock \\
pthread\_mutexattr\_t & Mutex attributes \\
pthread\_cond\_t & Condition variable \\
pthread\_condattr\_t & Condition attributes \\
pthread\_rwlock\_t & Read-write lock \\
pthread\_barrier\_t & Barrier sync \\
pthread\_spinlock\_t & Spinlock \\
pthread\_key\_t & Thread-local key \\
pthread\_once\_t & Once control \\
\end{tabular}
\end{tcolorbox}

% ============================================================================
% THREAD MANAGEMENT
% ============================================================================
\styledSection{\faCogs}{Thread Management}

\styledSubSection{Create \& Join}
\begin{lstlisting}
int pthread_create(
    pthread_t *thread,
    const pthread_attr_t *attr,
    void *(*start_routine)(void*),
    void *arg
);

int pthread_join(
    pthread_t thread,
    void **retval  // Get return value
);
\end{lstlisting}

\styledSubSection{Other Functions}
\begin{tcolorbox}[enhanced, colback=bgcard, colframe=bgcard, boxrule=0pt, arc=3pt, left=3pt, right=3pt, top=2pt, bottom=2pt]
\tablestyle
\begin{tabular}{@{}>{\ttfamily\color{primaryblue}}l@{\hspace{3pt}}>{\color{textmuted}}l@{}}
pthread\_self() & Get current thread ID \\
pthread\_equal(t1,t2) & Compare thread IDs \\
pthread\_detach(t) & Auto-cleanup on exit \\
pthread\_exit(ret) & Exit current thread \\
\end{tabular}
\end{tcolorbox}

\begin{tipbox}
\fontsize{9pt}{11pt}\selectfont
\textcolor{textlight}{\textbf{Detached} threads free resources automatically.}\\
\textcolor{textlight}{\textbf{Joinable} (default) must be joined to avoid memory leaks.}
\end{tipbox}

\styledSubSection{Basic Example}
\begin{lstlisting}
void* worker(void* arg) {
    int id = *(int*)arg;
    printf("Thread %d running\n", id);
    return NULL;
}

int main() {
    pthread_t t;
    int arg = 42;
    pthread_create(&t, NULL, worker, &arg);
    pthread_join(t, NULL);
    return 0;
}
\end{lstlisting}

% ============================================================================
% THREAD ATTRIBUTES
% ============================================================================
\styledSection{\faSlidersH}{Thread Attributes}

\begin{lstlisting}
pthread_attr_t attr;
pthread_attr_init(&attr);

// Detach state
pthread_attr_setdetachstate(&attr, 
    PTHREAD_CREATE_DETACHED);
    // or PTHREAD_CREATE_JOINABLE

// Stack size (e.g., 2MB)
pthread_attr_setstacksize(&attr, 
    2 * 1024 * 1024);

// Scheduling policy
pthread_attr_setschedpolicy(&attr,
    SCHED_FIFO);  // or SCHED_RR

pthread_attr_destroy(&attr);
\end{lstlisting}

% ============================================================================
% MUTEXES
% ============================================================================
\styledSection{\faLock}{Mutexes}

\styledSubSection{Basic Operations}
\begin{lstlisting}
// Static initialization
pthread_mutex_t m = PTHREAD_MUTEX_INITIALIZER;

// Dynamic initialization
pthread_mutex_init(&m, NULL);

// Lock operations
pthread_mutex_lock(&m);     // Blocking
pthread_mutex_trylock(&m);  // Non-blocking
pthread_mutex_unlock(&m);   // Release

// Cleanup
pthread_mutex_destroy(&m);
\end{lstlisting}

\styledSubSection{Timed Lock}
\begin{lstlisting}
struct timespec ts;
clock_gettime(CLOCK_REALTIME, &ts);
ts.tv_sec += 5;  // 5 second timeout

int rc = pthread_mutex_timedlock(&m, &ts);
if (rc == ETIMEDOUT) {
    // Handle timeout
}
\end{lstlisting}

\styledSubSection{Mutex Types}
\begin{tcolorbox}[enhanced, colback=bgcard, colframe=bgcard, boxrule=0pt, arc=3pt, left=3pt, right=3pt, top=2pt, bottom=2pt]
\tablestyle
\begin{tabular}{@{}>{\ttfamily\color{accentpurple}}l@{\hspace{3pt}}>{\color{textmuted}}l@{}}
NORMAL & Deadlock if re-locked \\
ERRORCHECK & Returns error on re-lock \\
RECURSIVE & Allows recursive locking \\
\end{tabular}
\end{tcolorbox}

\begin{lstlisting}
pthread_mutexattr_t attr;
pthread_mutexattr_init(&attr);
pthread_mutexattr_settype(&attr,
    PTHREAD_MUTEX_RECURSIVE);
pthread_mutex_init(&m, &attr);
pthread_mutexattr_destroy(&attr);
\end{lstlisting}

\begin{warnbox}
\fontsize{9pt}{11pt}\selectfont
\textcolor{textlight}{Always unlock in the \textbf{same thread} that locked!\\Never unlock a mutex you don't own.}
\end{warnbox}

% ============================================================================
% CONDITION VARIABLES
% ============================================================================
\styledSection{\faBell}{Condition Variables}

\styledSubSection{Basic Operations}
\begin{lstlisting}
// Static initialization
pthread_cond_t c = PTHREAD_COND_INITIALIZER;

// Wait (atomically releases mutex)
pthread_cond_wait(&c, &mutex);

// Wait with timeout
pthread_cond_timedwait(&c, &mutex, &ts);

// Signal one waiting thread
pthread_cond_signal(&c);

// Signal ALL waiting threads
pthread_cond_broadcast(&c);

// Cleanup
pthread_cond_destroy(&c);
\end{lstlisting}

\styledSubSection{Correct Usage Pattern}
\begin{lstlisting}
// === WAITER ===
pthread_mutex_lock(&mutex);
while (!condition) {    // Always WHILE!
    pthread_cond_wait(&cond, &mutex);
}
// ... do protected work ...
pthread_mutex_unlock(&mutex);

// === SIGNALER ===
pthread_mutex_lock(&mutex);
condition = true;
pthread_cond_signal(&cond);
pthread_mutex_unlock(&mutex);
\end{lstlisting}

\begin{warnbox}
\fontsize{9pt}{11pt}\selectfont
\textcolor{textlight}{Always use \texttt{while}, never \texttt{if}!\\Spurious wakeups can and do occur.}
\end{warnbox}

\styledSubSection{Signal vs Broadcast}
\begin{tcolorbox}[enhanced, colback=bgcard, colframe=bgcard, boxrule=0pt, arc=3pt, left=3pt, right=3pt, top=2pt, bottom=2pt]
\tablestyle
\begin{tabular}{@{}>{\ttfamily\color{primaryblue}}l@{\hspace{4pt}}>{\color{textmuted}}l@{}}
signal & Wake \textbf{one} waiting thread \\
broadcast & Wake \textbf{all} waiting threads \\
\end{tabular}
\end{tcolorbox}

\textcolor{textmuted}{\footnotesize Use \texttt{broadcast} for: shutdown signals, state changes affecting multiple threads.}

% ============================================================================
% READ-WRITE LOCKS
% ============================================================================
\styledSection{\faBook}{Read-Write Locks}

\textcolor{textmuted}{\footnotesize Multiple concurrent readers \textbf{OR} single exclusive writer.}

\begin{lstlisting}
pthread_rwlock_t rw = 
    PTHREAD_RWLOCK_INITIALIZER;

// Read lock (shared access)
pthread_rwlock_rdlock(&rw);
pthread_rwlock_tryrdlock(&rw);

// Write lock (exclusive access)
pthread_rwlock_wrlock(&rw);
pthread_rwlock_trywrlock(&rw);

// Unlock (same for both)
pthread_rwlock_unlock(&rw);

pthread_rwlock_destroy(&rw);
\end{lstlisting}

\begin{tipbox}
\fontsize{9pt}{11pt}\selectfont
\textcolor{textlight}{Best when reads >> writes. For balanced read/write, use a simple mutex instead.}
\end{tipbox}

% ============================================================================
% BARRIERS
% ============================================================================
\styledSection{\faUsers}{Barriers}

\textcolor{textmuted}{\footnotesize Synchronize exactly N threads at a rendezvous point.}

\begin{lstlisting}
pthread_barrier_t b;

// Initialize for N threads
pthread_barrier_init(&b, NULL, N);

// Wait (blocks until N threads arrive)
int rc = pthread_barrier_wait(&b);
if (rc == PTHREAD_BARRIER_SERIAL_THREAD) {
    // Exactly ONE thread gets this
    // Use for cleanup, etc.
}

pthread_barrier_destroy(&b);
\end{lstlisting}

% ============================================================================
% SPINLOCKS
% ============================================================================
\styledSection{\faSync}{Spinlocks}

\textcolor{textmuted}{\footnotesize Busy-wait lock for \textit{extremely short} critical sections.}

\begin{lstlisting}
pthread_spinlock_t s;

pthread_spin_init(&s, 
    PTHREAD_PROCESS_PRIVATE);

pthread_spin_lock(&s);    // Busy waits!
pthread_spin_trylock(&s); // Non-blocking
pthread_spin_unlock(&s);

pthread_spin_destroy(&s);
\end{lstlisting}

\begin{warnbox}
\fontsize{9pt}{11pt}\selectfont
\textcolor{textlight}{Burns CPU cycles while waiting! Only use for lock hold times < context switch time.}
\end{warnbox}

% ============================================================================
% THREAD-LOCAL STORAGE
% ============================================================================
\styledSection{\faKey}{Thread-Local Storage}

\styledSubSection{pthread\_key API}
\begin{lstlisting}
pthread_key_t key;

// Create key with destructor (once)
pthread_key_create(&key, destructor_fn);

// Set value for THIS thread only
pthread_setspecific(key, value);

// Get value for THIS thread
void* val = pthread_getspecific(key);

pthread_key_delete(key);
\end{lstlisting}

\styledSubSection{\_\_thread Keyword (GCC/Clang)}
\begin{lstlisting}
__thread int counter = 0;
__thread void* my_data = NULL;

// Each thread has its own copy!
// Simpler than pthread_key API
\end{lstlisting}

% ============================================================================
% ONCE INITIALIZATION
% ============================================================================
\styledSection{\faCheckDouble}{Once Initialization}

\textcolor{textmuted}{\footnotesize Execute initialization code exactly once, thread-safely.}

\begin{lstlisting}
pthread_once_t once = PTHREAD_ONCE_INIT;

void init_function(void) {
    // This runs exactly ONCE
    // regardless of how many threads call
}

// Called from any/all threads:
pthread_once(&once, init_function);
\end{lstlisting}

% ============================================================================
% CANCELLATION
% ============================================================================
\styledSection{\faBan}{Thread Cancellation}

\begin{lstlisting}
// Request thread cancellation
pthread_cancel(thread);

// Enable/disable cancellation
pthread_setcancelstate(
    PTHREAD_CANCEL_ENABLE, NULL);
    // or PTHREAD_CANCEL_DISABLE

// Set cancellation type
pthread_setcanceltype(
    PTHREAD_CANCEL_DEFERRED, NULL);
    // or PTHREAD_CANCEL_ASYNCHRONOUS

// Create explicit cancellation point
pthread_testcancel();
\end{lstlisting}

\styledSubSection{Cleanup Handlers}
\begin{lstlisting}
pthread_cleanup_push(cleanup_fn, arg);
// ... cancellable code ...
pthread_cleanup_pop(1);  // 1 = execute handler
\end{lstlisting}

% ============================================================================
% SEMAPHORES
% ============================================================================
\styledSection{\faTrafficLight}{Semaphores}

\textcolor{accentorange}{\texttt{\#include <semaphore.h>}}

\begin{lstlisting}
sem_t sem;

// Initialize with count
sem_init(&sem, 0, initial_value);

// Wait (decrement, blocks if 0)
sem_wait(&sem);
sem_trywait(&sem);
sem_timedwait(&sem, &ts);

// Post (increment, wake waiters)
sem_post(&sem);

// Get current value
sem_getvalue(&sem, &val);

sem_destroy(&sem);
\end{lstlisting}

\begin{tipbox}
\fontsize{9pt}{11pt}\selectfont
\textcolor{textlight}{\textbf{Mutex} = binary semaphore (0 or 1)}\\
\textcolor{textlight}{\textbf{Semaphore} = counting (0 to N)}
\end{tipbox}

% ============================================================================
% ERROR CODES
% ============================================================================
\styledSection{\faExclamationCircle}{Error Codes}

\begin{tcolorbox}[enhanced, colback=bgcard, colframe=bgcard, boxrule=0pt, arc=3pt, left=3pt, right=3pt, top=2pt, bottom=2pt]
\tablestyle
\begin{tabular}{@{}>{\ttfamily\color{codered}}l@{\hspace{4pt}}>{\color{textmuted}}l@{}}
EAGAIN & Insufficient resources \\
EBUSY & Resource is busy \\
EDEADLK & Deadlock detected \\
EINVAL & Invalid argument \\
ENOMEM & Out of memory \\
EPERM & Operation not permitted \\
ESRCH & No such thread \\
ETIMEDOUT & Operation timed out \\
\end{tabular}
\end{tcolorbox}

\begin{lstlisting}
int rc = pthread_create(&t, NULL, fn, arg);
if (rc != 0) {
    fprintf(stderr, "Error: %s\n", 
            strerror(rc));
    exit(EXIT_FAILURE);
}
\end{lstlisting}

% ============================================================================
% SCHEDULING
% ============================================================================
\styledSection{\faClock}{Scheduling}

\begin{lstlisting}
// Get/Set scheduling parameters
pthread_getschedparam(t, &policy, &param);
pthread_setschedparam(t, policy, &param);

// Set thread priority
pthread_setschedprio(t, priority);

// Yield CPU to other threads
sched_yield();

// Get priority range for policy
int min = sched_get_priority_min(policy);
int max = sched_get_priority_max(policy);
\end{lstlisting}

\styledSubSection{Scheduling Policies}
\begin{tcolorbox}[enhanced, colback=bgcard, colframe=bgcard, boxrule=0pt, arc=3pt, left=3pt, right=3pt, top=2pt, bottom=2pt]
\tablestyle
\begin{tabular}{@{}>{\ttfamily\color{accentpurple}}l@{\hspace{3pt}}>{\color{textmuted}}l@{}}
SCHED\_OTHER & Default time-sharing \\
SCHED\_FIFO & Real-time FIFO \\
SCHED\_RR & Real-time round-robin \\
\end{tabular}
\end{tcolorbox}

% ============================================================================
% COMMON PATTERNS
% ============================================================================
\styledSection{\faLightbulb}{Common Patterns}

\styledSubSection{Producer-Consumer Queue}
\begin{lstlisting}
// === PRODUCER ===
pthread_mutex_lock(&mutex);
while (queue_full)
    pthread_cond_wait(&not_full, &mutex);
enqueue(item);
pthread_cond_signal(&not_empty);
pthread_mutex_unlock(&mutex);

// === CONSUMER ===
pthread_mutex_lock(&mutex);
while (queue_empty)
    pthread_cond_wait(&not_empty, &mutex);
item = dequeue();
pthread_cond_signal(&not_full);
pthread_mutex_unlock(&mutex);
\end{lstlisting}

\styledSubSection{Thread Pool Worker}
\begin{lstlisting}
void* worker(void* pool_ptr) {
    while (1) {
        pthread_mutex_lock(&mutex);
        while (queue_empty && !shutdown)
            pthread_cond_wait(&job_ready, &mutex);
        if (shutdown && queue_empty) {
            pthread_mutex_unlock(&mutex);
            break;
        }
        job = dequeue_job();
        pthread_mutex_unlock(&mutex);
        execute_job(job);
    }
    return NULL;
}
\end{lstlisting}

\styledSubSection{Graceful Shutdown}
\begin{lstlisting}
// Signal all workers to stop
pthread_mutex_lock(&mutex);
shutdown = true;
pthread_cond_broadcast(&cond);  // Wake ALL
pthread_mutex_unlock(&mutex);

// Wait for all threads to finish
for (int i = 0; i < num_threads; i++)
    pthread_join(threads[i], NULL);

// Clean up resources
pthread_mutex_destroy(&mutex);
pthread_cond_destroy(&cond);
\end{lstlisting}

% ============================================================================
% GOLDEN RULES
% ============================================================================
\styledSection{\faStar}{Golden Rules}

\begin{infobox}
\fontsize{9pt}{11pt}\selectfont
\begin{enumerate}[leftmargin=12pt, itemsep=1pt]
\item \textcolor{successgreen}{\textbf{Always}} use \texttt{while} with \texttt{cond\_wait}
\item \textcolor{successgreen}{\textbf{Lock}} before checking/modifying shared state
\item \textcolor{successgreen}{\textbf{Unlock}} what you lock (same thread!)
\item \textcolor{successgreen}{\textbf{Join}} or \textbf{detach} every thread created
\item \textcolor{successgreen}{\textbf{Destroy}} every resource you initialize
\item \textcolor{dangerred}{\textbf{Don't}} pass stack variables without sync
\item \textcolor{successgreen}{\textbf{Lock in consistent order}} to avoid deadlock
\item \textcolor{successgreen}{\textbf{Minimize}} critical section duration
\end{enumerate}
\end{infobox}

% ============================================================================
% COMPILATION
% ============================================================================
\styledSection{\faTerminal}{Compilation}

\begin{lstlisting}[language=bash]
# Basic compilation (Linux/macOS)
gcc -pthread program.c -o program

# With debugging and thread sanitizer
gcc -pthread -g -fsanitize=thread \
    program.c -o program

# With additional libraries
gcc -pthread program.c -o program -lm -lrt
\end{lstlisting}

% ============================================================================
% QUICK REFERENCE TABLE
% ============================================================================
\styledSection{\faTable}{Quick Reference}

\begin{tcolorbox}[enhanced, colback=bgcard, colframe=primaryblue!50, boxrule=0.5pt, arc=3pt, left=3pt, right=3pt, top=2pt, bottom=2pt]
\fontsize{9pt}{11pt}\selectfont
\renewcommand{\arraystretch}{1.3}
\tablestyle
\begin{tabular}{@{}>{\color{textlight}}l@{\hspace{3pt}}>{\ttfamily\color{primaryblue}}l@{}}
\rowcolor{primaryblue!20!bgdark} \multicolumn{2}{@{}l}{\textbf{\color{white}Thread Operations}} \\
Create thread & pthread\_create \\
Wait for thread & pthread\_join \\
Exit thread & pthread\_exit \\
\rowcolor{primaryblue!20!bgdark} \multicolumn{2}{@{}l}{\textbf{\color{white}Mutex Operations}} \\
Lock (blocking) & pthread\_mutex\_lock \\
Lock (try) & pthread\_mutex\_trylock \\
Unlock & pthread\_mutex\_unlock \\
\rowcolor{primaryblue!20!bgdark} \multicolumn{2}{@{}l}{\textbf{\color{white}Condition Variables}} \\
Wait for signal & pthread\_cond\_wait \\
Signal one & pthread\_cond\_signal \\
Signal all & pthread\_cond\_broadcast \\
\rowcolor{primaryblue!20!bgdark} \multicolumn{2}{@{}l}{\textbf{\color{white}RW Locks \& Barriers}} \\
Read lock & pthread\_rwlock\_rdlock \\
Write lock & pthread\_rwlock\_wrlock \\
Barrier wait & pthread\_barrier\_wait \\
\end{tabular}
\end{tcolorbox}

\vfill

% Footer
\begin{center}
\begin{tikzpicture}
\draw[primaryblue!50, line width=0.5pt] (0,0) -- (\linewidth,0);
\end{tikzpicture}
\vspace{1pt}

{\fontsize{8pt}{9pt}\selectfont\color{textmuted}
\texttt{\color{accentorange}\#include <pthread.h>} \quad\textbar\quad 
Link: \texttt{\color{primaryblue}-pthread} \quad\textbar\quad 
POSIX 1003.1c-1995
}
\end{center}

\end{multicols}

\end{document}
